\documentclass[12pt,a4paper]{article}

\usepackage[utf8]{inputenc}
\usepackage[T1]{fontenc}
\usepackage[brazil]{babel}
\usepackage{mathptmx}
\usepackage{geometry}
\usepackage{graphicx}
\usepackage{booktabs}
\usepackage{array}
\usepackage{indentfirst}
\usepackage{url}
\usepackage{microtype}

\geometry{
  a4paper,
  left=2.5cm,
  right=2.5cm,
  top=2.5cm,
  bottom=2.5cm
}

\pagestyle{empty}
\setlength{\parindent}{0.75cm}
\setlength{\parskip}{0pt}
\renewcommand{\baselinestretch}{1.0}
\emergencystretch=2em

\newcommand{\linhaautor}[3]{\noindent\textbf{#1}$^{#2}$ - #3\par}
\newcommand{\linhaafiliacao}[2]{%
  \noindent\hangindent=0.38cm\hangafter=1
  $^{#1}$ #2\par
}

\begin{document}

\vspace*{2pt}

\begin{center}
\begin{minipage}{0.96\textwidth}
\centering
\fontsize{12.2}{13.6}\selectfont\bfseries
COMPARAÇÃO DE MODELOS DE APRENDIZADO DE MÁQUINA\\
E REDES AVANÇADAS NA PREVISÃO DIÁRIA E MENSAL\\
DE IRRADIÂNCIA GLOBAL HORIZONTAL EM LOCALIDADES\\
DE FÁBRICAS DE VEÍCULOS ELÉTRICOS%
\end{minipage}
\end{center}

\vspace{24pt}

\begingroup
\fontsize{11.6}{13.4}\selectfont
\linhaautor{Gustavo De Melo Nascimento}{1}{g239679@dac.unicamp.br}
\linhaautor{Thales Wulfert Cabral}{2}{t264377@dac.unicamp.br}
\linhaautor{Gustavo Fraidenraich}{3}{gfraiden@unicamp.br}

\vspace{9pt}

\linhaafiliacao{1}{Aluno, Faculdade de Engenharia Elétrica e de Computação (FEEC), Universidade Estadual de Campinas (UNICAMP) - Campinas, SP, Brasil}
\linhaafiliacao{2}{Coorientador, Faculdade de Engenharia Elétrica e de Computação (FEEC), Universidade Estadual de Campinas (UNICAMP) - Campinas, SP, Brasil}
\linhaafiliacao{3}{Orientador, DECOM, Faculdade de Engenharia Elétrica e de Computação (FEEC), Universidade Estadual de Campinas (UNICAMP) - Campinas, SP, Brasil}
\endgroup

\vspace{26pt}

\begingroup
\fontsize{11.6}{13.5}\selectfont
\noindent\textbf{\textit{Resumo.}} \textit{Este artigo apresenta um fluxo computacional para previsão da Irradiância Global Horizontal (GHI) em dez localidades associadas a fábricas de veículos elétricos. Foram utilizados dados horários oficiais da NSRDB, disponibilizada pelo National Renewable Energy Laboratory (NREL), no produto GOES Aggregated PSM v4, para o período de 2019 a 2024. A série horária foi convertida para médias diárias e médias mensais. Após validação, quantização em 128 níveis, normalização e criação de variáveis temporais, foram avaliados sete modelos sob o mesmo protocolo experimental: XGBoost, MLP, RNN, LSTM, DilatedRNN, DeepAR experimental e DeepNPTS aproximado. A divisão entre treino e teste preservou a ordem temporal, com 80\% dos exemplos para treinamento e 20\% para avaliação. Entre os modelos de referência, a média dos melhores resultados por localidade apresentou, na escala diária, MAE de 37,77 W/m$^2$, RMSE de 49,72 W/m$^2$, R$^2$ de 0,66 e nRMSE de 26,86\%. Na escala mensal, a média correspondente foi MAE de 13,63 W/m$^2$, RMSE de 16,77 W/m$^2$, R$^2$ de 0,89 e nRMSE de 8,28\%. Na comparação conjunta, a DilatedRNN foi competitiva na escala diária e o DeepNPTS aproximado apresentou o menor RMSE médio na escala mensal.}

\vspace{12pt}

\noindent\textbf{\textit{Palavras-chave:}} Irradiância solar, Séries temporais, XGBoost, MLP, RNN, LSTM, DilatedRNN, DeepAR, DeepNPTS
\endgroup

\vspace{24pt}

\noindent\textbf{1. INTRODUÇÃO}

A previsão de irradiância solar é uma etapa relevante para estudos de planejamento energético, operação de sistemas fotovoltaicos e avaliação de recursos solares. A Irradiância Global Horizontal (GHI) representa a radiação solar incidente sobre uma superfície horizontal e é uma das variáveis mais utilizadas para caracterizar a disponibilidade de energia solar em uma localidade. Em aplicações fotovoltaicas, a análise de séries históricas de GHI apoia estudos preliminares de viabilidade, comparação entre regiões e dimensionamento energético (Pinho e Galdino, 2014; Villalva e Gazoli, 2012).

Com o crescimento da eletrificação do setor de transportes, localidades associadas a fábricas de veículos elétricos tornam-se pontos de interesse para análises energéticas. Essas instalações combinam alto consumo de energia, cadeias produtivas estratégicas e possibilidade de integração com geração solar fotovoltaica local. Neste trabalho, as fábricas são utilizadas como pontos geográficos de análise, sem assumir que elas alterem ou causem padrões específicos de irradiância.

Séries medidas em estações meteorológicas podem apresentar limitações de cobertura espacial, períodos curtos de medição e falhas de registro. Bases satelitais, como o NSRDB, permitem obter séries padronizadas para diferentes coordenadas geográficas, tornando possível comparar localidades distintas sob uma mesma metodologia. Esse tipo de abordagem é especialmente útil em estudos computacionais que envolvem múltiplas regiões e requerem reprodutibilidade.

O objetivo deste artigo é prever a GHI média do próximo período a partir do histórico observado até o período atual. O fluxo diário utiliza médias diárias para prever o dia seguinte, enquanto o fluxo mensal utiliza médias mensais para prever o mês seguinte. Para isso, foi desenvolvido um fluxo de aprendizado de máquina baseado em dados oficiais da NSRDB/NREL, com validação, transformação temporal e comparação entre sete modelos de regressão e séries temporais: XGBoost, MLP, RNN, LSTM, DilatedRNN, DeepAR experimental e DeepNPTS aproximado. A principal contribuição está na construção de um procedimento reprodutível e auditável para avaliar modelos de previsão de curto prazo em dez localidades distribuídas no Brasil, Estados Unidos e México, mantendo a mesma preparação de dados, o mesmo corte temporal e as mesmas métricas para todos os modelos.

\vspace{12pt}

\noindent\textbf{2. METODOLOGIA}

\vspace{12pt}

\noindent\textbf{2.1 Dados utilizados}

Os dados utilizados foram obtidos a partir da National Solar Radiation Database (NSRDB), base amplamente utilizada para séries históricas de radiação solar (Wilcox, 2012; Sengupta et al., 2018), disponibilizada pelo National Renewable Energy Laboratory (NREL). Utilizou-se o produto GOES Aggregated PSM v4, com a variável GHI em resolução horária de 60 minutos e unidade W/m$^2$. O período analisado compreende 1 de janeiro de 2019 a 31 de dezembro de 2024, totalizando 2.192 dias por localidade. A aquisição e a organização dos dados foram realizadas em Python, com apoio de bibliotecas para processamento de dados solares, como o pvlib (PVLIB Python Development Team, 2026).

Foram avaliadas dez localidades associadas a fábricas de veículos elétricos: BYD Camaçari, Tesla Gigafactory Nevada, Tesla Gigafactory Texas, Hyundai Metaplant Georgia, Rivian Normal, Tesla Fremont Factory, Lucid AMP 1 Casa Grande, GM Factory Zero, Ford Rouge Electric Vehicle Center e BMW San Luis Potosí. A identificação das fábricas foi documentada a partir de fontes institucionais, e as coordenadas foram associadas ao centro do elemento industrial correspondente no OpenStreetMap, consultado por meio do Nominatim. Cada localidade foi tratada de forma independente, com um conjunto de dados, modelos e métricas próprios.

\vspace{12pt}

\noindent\textbf{2.2 Pré-processamento}

Inicialmente, os registros horários de GHI foram agregados por média diária. Essa transformação preserva a unidade W/m$^2$, pois representa a média da potência solar incidente ao longo do dia, e não a energia acumulada diária. Em um segundo fluxo, as médias diárias foram agregadas por mês civil, gerando uma série mensal de GHI médio. Assim, o estudo avalia dois horizontes de curto prazo: previsão do dia seguinte e previsão do mês seguinte.

Em seguida, a série foi validada quanto à cobertura temporal, datas duplicadas, valores ausentes, valores negativos, unidade, origem dos dados, período e compatibilidade geográfica. Também foram preservados metadados de proveniência, como coordenadas da fábrica, ponto da grade solar retornado pela API, produto, endpoint, unidade, intervalo temporal, agregação e data de coleta. Essa etapa permite auditar a origem dos dados e reduz o risco de utilizar arquivos sintéticos ou incompletos na modelagem.

Após a limpeza, a série de GHI de cada fluxo foi quantizada em 128 níveis inteiros, de 0 a 127, e posteriormente normalizada para o intervalo entre 0 e 1. Os parâmetros de escala foram ajustados apenas no conjunto de treinamento e reaplicados ao conjunto de teste, reduzindo o risco de vazamento de informação futura.

\vspace{12pt}

\noindent\textbf{2.3 Variáveis temporais e alvo}

O problema foi formulado como regressão de série temporal. No fluxo diário, para uma linha associada ao dia $t$, o alvo é a GHI do dia $t+1$. As variáveis de entrada foram construídas apenas com informações disponíveis até o dia $t$, incluindo quatro defasagens e três médias móveis. A notação das defasagens é relativa ao alvo: GHI$_{t-1}$ corresponde ao valor observado no período imediatamente anterior ao previsto.

\vspace{6pt}

\begin{center}
\begin{tabular}{ll}
\toprule
\textbf{Variável} & \textbf{Descrição} \\
\midrule
GHI$_{t-1}$ & GHI observada no dia $t$ \\
GHI$_{t-2}$ & GHI observada no dia $t-1$ \\
GHI$_{t-3}$ & GHI observada no dia $t-2$ \\
GHI$_{t-7}$ & GHI observada no dia $t-6$ \\
Média móvel 3d & Média dos últimos 3 dias até $t$ \\
Média móvel 7d & Média dos últimos 7 dias até $t$ \\
Média móvel 30d & Média dos últimos 30 dias até $t$ \\
\bottomrule
\end{tabular}
\end{center}

\vspace{6pt}

A janela de 30 dias remove as primeiras 29 linhas da série, e a ausência de alvo para o último dia remove mais uma linha. Assim, a base final de modelagem contém 2.162 exemplos por localidade.

No fluxo mensal, a mesma lógica foi aplicada sobre a série agregada por mês. Nesse caso, o alvo é a GHI média do mês seguinte, com defasagens mensais de 1, 2, 3 e 6 meses e médias móveis de 3, 6 e 12 meses. A base mensal contém 60 exemplos por localidade após a remoção das linhas sem histórico completo e sem alvo futuro.

\vspace{12pt}

\noindent\textbf{2.4 Modelos de referência}

O fluxo computacional foi estruturado inicialmente com quatro modelos de referência, usando exatamente as mesmas entradas, o mesmo alvo e a mesma divisão temporal. O XGBoost utiliza árvores de decisão combinadas por gradiente e apresenta bom desempenho em problemas tabulares não lineares (Chen e Guestrin, 2016). A MLP recebeu as sete variáveis de entrada de cada fluxo e foi configurada como uma rede neural totalmente conectada. Além disso, foram incorporadas arquiteturas recorrentes RNN e LSTM, nas quais as sete variáveis temporais são reinterpretadas como uma sequência curta de sete passos, com uma variável por passo. Redes neurais desse tipo são capazes de aproximar relações não lineares entre entradas históricas e o alvo previsto (Haykin, 2009; Géron, 2022). A formulação seguiu uma estratégia supervisionada de previsão temporal, na qual observações passadas são usadas para estimar um valor futuro (Hyndman e Athanasopoulos, 2021).

\vspace{6pt}

\begin{center}
\small
\begin{tabular}{p{2.2cm}p{5.0cm}p{6.6cm}}
\toprule
\textbf{Modelo} & \textbf{Entrada} & \textbf{Configuração utilizada} \\
\midrule
XGBoost & Sete features temporais tabulares & 300 estimadores, profundidade máxima 3, taxa de aprendizagem 0,05, \textit{subsample} 0,9, \textit{colsample} 0,9 e objetivo \textit{squarederror}. \\
MLP & Sete features temporais tabulares & Duas camadas ocultas, com 64 e 32 neurônios, ativação ReLU, otimizador Adam, taxa inicial 0,001 e máximo de 1000 iterações. \\
RNN & Sequência curta de sete passos & SimpleRNN com 32 unidades, camada densa com 16 neurônios ReLU e saída sigmoidal para manter a previsão em $[0,1]$. \\
LSTM & Sequência curta de sete passos & LSTM com 32 unidades, camada densa com 16 neurônios ReLU e saída sigmoidal para manter a previsão em $[0,1]$. \\
\bottomrule
\end{tabular}
\end{center}

\vspace{6pt}

Nos modelos recorrentes, a matriz de entrada originalmente tabular foi reorganizada no formato ``amostras $\times$ passos $\times$ variável'', resultando em sete passos temporais e uma variável por passo. O treinamento das redes Keras utilizou perda MSE, otimizador Adam com taxa de aprendizagem 0,001, até 80 épocas, lote de 32 amostras, validação interna de 15\% do conjunto de treinamento e parada antecipada com paciência de 10 épocas quando aplicável. As previsões foram limitadas ao intervalo $[0,1]$ antes da conversão para W/m$^2$.

\vspace{12pt}

\noindent\textbf{2.5 Avaliação}

A divisão entre treinamento e teste foi feita de forma cronológica, sem embaralhamento. Foram utilizados 80\% dos exemplos para treinamento e 20\% para teste em cada fluxo. Na escala diária, o período de teste corresponde aos alvos entre 26 de outubro de 2023 e 31 de dezembro de 2024. Na escala mensal, o teste corresponde aos alvos mensais de 2024. Assim, os experimentos avaliam previsão de curto prazo em período não visto durante o treinamento, sem produzir projeções para anos posteriores à base observada.

As métricas calculadas foram erro absoluto médio (MAE), erro quadrático médio (MSE), raiz do erro quadrático médio (RMSE), coeficiente de determinação (R$^2$) e RMSE normalizado (nRMSE), calculado em relação à média da GHI observada no teste. Como os modelos aprendem na escala quantizada e normalizada, as métricas foram salvas em duas escalas: a escala normalizada, usada internamente no treinamento, e a escala física aproximada em W/m$^2$. A desnormalização foi feita com os parâmetros ajustados no treinamento de cada localidade:

\vspace{6pt}

\begin{center}
$\widehat{GHI}_{W/m^2} = \widehat{GHI}_{norm} \cdot (GHI_{max,treino} - GHI_{min,treino}) + GHI_{min,treino}$.
\end{center}

\vspace{6pt}

O critério principal de comparação foi o maior R$^2$ na escala em W/m$^2$, pois essa versão está alinhada à interpretação física dos erros. O nRMSE percentual foi calculado como $100 \cdot RMSE_{W/m^2} / \overline{GHI}_{real,W/m^2}$. O pipeline também registra média, desvio-padrão e COV horário ($\sigma/\mu$) quando os dados horários originais estão disponíveis antes da agregação diária. Nos CSVs diários já agregados, essa variabilidade intradiária não pode ser reconstruída, e a estatística é marcada como indisponível para COV horário.

\vspace{12pt}

\noindent\textbf{2.6 Modelos avançados avaliados}

Para ampliar a comparação, foram avaliadas três abordagens avançadas sob o mesmo protocolo dos modelos de referência: DilatedRNN, DeepAR\_exp e DeepNPTS\_aprox. Esses modelos utilizaram a mesma base processada, as mesmas variáveis temporais, o mesmo alvo, a mesma divisão cronológica e as mesmas métricas. Dessa forma, a comparação entre modelos de referência e modelos avançados é feita em condições equivalentes. As arquiteturas foram inspiradas em propostas de redes recorrentes dilatadas, previsão probabilística autoregressiva e bibliotecas de modelagem probabilística de séries temporais (Chang et al., 2017; Salinas et al., 2017; Alexandrov et al., 2019).

\vspace{6pt}

\begin{center}
\small
\begin{tabular}{p{2.8cm}p{10.0cm}}
\toprule
\textbf{Modelo} & \textbf{Descrição da abordagem} \\
\midrule
DilatedRNN & Rede recorrente com ramos dilatados, usando fatores de dilatação 1, 2 e 4 sobre a sequência curta de sete passos. A proposta foi capturar dependências temporais em diferentes espaçamentos dentro da mesma janela de entrada, seguindo a ideia de conexões recorrentes dilatadas. \\
DeepAR\_exp & Modelo recorrente experimental inspirado no DeepAR, com estrutura LSTM e perda probabilística gaussiana. A saída estima média e dispersão, mas a previsão pontual usada nas métricas corresponde à média prevista. \\
DeepNPTS\_aprox & Aproximação não paramétrica inspirada em NPTS, baseada em vizinhos históricos ponderados por similaridade e recência. O objetivo foi testar uma alternativa menos paramétrica para séries suavizadas, mantendo o mesmo protocolo de avaliação dos demais modelos. \\
\bottomrule
\end{tabular}
\end{center}

\vspace{6pt}

No artigo, esses três modelos são analisados em conjunto com os demais. A distinção entre modelos de referência e modelos avançados serve apenas para organizar a discussão metodológica, já que todos compartilham a mesma preparação dos dados e o mesmo critério de avaliação.

\vspace{12pt}

\noindent\textbf{3. RESULTADOS E DISCUSSÃO}

Os modelos de referência foram capazes de capturar parte relevante da dinâmica temporal da GHI diária. Considerando o melhor modelo de referência em cada localidade, a média geral foi de MAE igual a 37,77 W/m$^2$, RMSE igual a 49,72 W/m$^2$, R$^2$ igual a 0,66 e nRMSE igual a 26,86\%. A RNN obteve o melhor desempenho em cinco das dez localidades, o XGBoost em três e a MLP em duas. A LSTM foi treinada e avaliada em todas as localidades, mas não apresentou o maior R$^2$ em nenhuma delas nesta configuração diária.

\vspace{6pt}

\begin{center}
\small
\begin{tabular}{lrrrr}
\toprule
\textbf{Modelo} & \textbf{MAE} & \textbf{RMSE} & \textbf{R$^2$} & \textbf{nRMSE} \\
 & \textbf{W/m$^2$} & \textbf{W/m$^2$} & & \textbf{\%} \\
\midrule
XGBoost & 38,42 & 50,77 & 0,64 & 27,41 \\
MLP & 37,85 & 50,08 & 0,65 & 27,05 \\
RNN & 37,90 & 49,91 & 0,65 & 26,95 \\
LSTM & 40,27 & 53,37 & 0,61 & 28,71 \\
\bottomrule
\end{tabular}
\end{center}

\vspace{6pt}

A tabela acima mostra a média das métricas diárias considerando todas as localidades para cada modelo. A RNN obteve o menor RMSE médio e o maior R$^2$ médio, com desempenho muito próximo ao da MLP. O XGBoost apresentou resultados competitivos, enquanto a LSTM teve maior erro médio nesta configuração. Como as diferenças entre XGBoost, MLP e RNN foram pequenas, a análise por localidade é necessária para evitar a interpretação de um vencedor universal.

\vspace{6pt}

\begin{center}
\small
\begin{tabular}{p{5.2cm}p{1.6cm}p{1.9cm}rrrr}
\toprule
\textbf{Localidade} & \textbf{País} & \textbf{Modelo} & \textbf{MAE} & \textbf{RMSE} & \textbf{R$^2$} & \textbf{nRMSE} \\
 & & & \textbf{W/m$^2$} & \textbf{W/m$^2$} & & \textbf{\%} \\
\midrule
BYD Camaçari & Brasil & MLP & 38,75 & 49,67 & 0,39 & 20,35 \\
Tesla Gigafactory Nevada & EUA & XGBoost & 28,47 & 39,26 & 0,86 & 19,13 \\
Tesla Gigafactory Texas & EUA & RNN & 42,99 & 54,92 & 0,57 & 28,02 \\
Hyundai Metaplant Georgia & EUA & RNN & 45,87 & 56,93 & 0,52 & 30,50 \\
Rivian Normal & EUA & MLP & 46,85 & 59,50 & 0,64 & 36,11 \\
Tesla Fremont Factory & EUA & XGBoost & 25,05 & 37,11 & 0,87 & 18,14 \\
Lucid AMP 1 Casa Grande & EUA & RNN & 23,99 & 36,47 & 0,82 & 15,70 \\
GM Factory Zero & EUA & RNN & 45,97 & 59,51 & 0,64 & 41,04 \\
Ford Rouge Electric Vehicle Center & EUA & XGBoost & 46,01 & 58,51 & 0,65 & 40,31 \\
BMW San Luis Potosí & México & RNN & 33,71 & 45,35 & 0,60 & 19,30 \\
\bottomrule
\end{tabular}
\end{center}

\vspace{6pt}

O menor RMSE foi observado em Lucid AMP 1 Casa Grande, com RMSE de 36,47 W/m$^2$ e R$^2$ de 0,82 usando RNN. Os maiores valores de R$^2$ ocorreram em Tesla Fremont Factory e Tesla Gigafactory Nevada, ambas com melhor desempenho por XGBoost e R$^2$ acima de 0,85. Esses casos indicam que localidades com padrões de irradiância mais regulares tendem a favorecer previsões de curto prazo baseadas em defasagens e médias móveis.

As maiores dificuldades ocorreram em localidades como GM Factory Zero e Ford Rouge Electric Vehicle Center, que apresentaram nRMSE superior a 40\%. Esse comportamento sugere maior variabilidade temporal da irradiância diária ou menor capacidade das variáveis históricas simples em representar mudanças meteorológicas rápidas. Mesmo nesses casos, os modelos mantiveram R$^2$ acima de 0,63, indicando que parte significativa da variabilidade ainda foi explicada.

No fluxo mensal, o desempenho foi superior ao observado na escala diária, como esperado, pois a agregação mensal suaviza variações meteorológicas rápidas. Considerando o melhor modelo de cada localidade, a média geral foi de MAE igual a 13,63 W/m$^2$, RMSE igual a 16,77 W/m$^2$, R$^2$ igual a 0,89 e nRMSE igual a 8,28\%. A RNN apresentou o melhor desempenho em quatro localidades, enquanto MLP e XGBoost foram superiores em três localidades cada.

\vspace{6pt}

\begin{center}
\small
\begin{tabular}{lrrrr}
\toprule
\textbf{Modelo} & \textbf{MAE} & \textbf{RMSE} & \textbf{R$^2$} & \textbf{nRMSE} \\
 & \textbf{W/m$^2$} & \textbf{W/m$^2$} & & \textbf{\%} \\
\midrule
XGBoost & 15,23 & 19,41 & 0,88 & 9,70 \\
MLP & 16,13 & 19,70 & 0,86 & 9,58 \\
RNN & 14,81 & 18,23 & 0,88 & 8,97 \\
LSTM & 20,60 & 24,69 & 0,78 & 11,99 \\
\bottomrule
\end{tabular}
\end{center}

\vspace{6pt}

Na média mensal por modelo, a RNN apresentou o menor MAE, menor RMSE e menor nRMSE, enquanto XGBoost obteve R$^2$ médio muito próximo. A LSTM novamente ficou abaixo dos demais modelos de referência. A melhora em relação ao fluxo diário não significa que o problema mensal seja equivalente ao diário; ela reflete a suavização da série quando valores diários são agregados por mês civil.

\vspace{6pt}

\begin{center}
\small
\begin{tabular}{p{5.2cm}p{1.6cm}p{1.9cm}rrrr}
\toprule
\textbf{Localidade} & \textbf{País} & \textbf{Modelo} & \textbf{MAE} & \textbf{RMSE} & \textbf{R$^2$} & \textbf{nRMSE} \\
 & & & \textbf{W/m$^2$} & \textbf{W/m$^2$} & & \textbf{\%} \\
\midrule
BYD Camaçari & Brasil & RNN & 16,74 & 22,03 & 0,65 & 9,33 \\
Tesla Gigafactory Nevada & EUA & MLP & 10,02 & 12,29 & 0,98 & 5,53 \\
Tesla Gigafactory Texas & EUA & MLP & 15,13 & 18,44 & 0,90 & 8,83 \\
Hyundai Metaplant Georgia & EUA & XGBoost & 12,05 & 15,98 & 0,92 & 8,06 \\
Rivian Normal & EUA & MLP & 20,45 & 23,73 & 0,90 & 13,19 \\
Tesla Fremont Factory & EUA & RNN & 10,13 & 12,17 & 0,98 & 5,52 \\
Lucid AMP 1 Casa Grande & EUA & XGBoost & 5,61 & 6,65 & 0,99 & 2,70 \\
GM Factory Zero & EUA & RNN & 13,04 & 16,12 & 0,96 & 10,06 \\
Ford Rouge Electric Vehicle Center & EUA & RNN & 12,10 & 15,38 & 0,96 & 9,59 \\
BMW San Luis Potosí & México & XGBoost & 21,07 & 24,91 & 0,69 & 10,05 \\
\bottomrule
\end{tabular}
\end{center}

\vspace{6pt}

A comparação entre os modelos mostra que, na escala diária, a RNN obteve o melhor resultado no maior número de localidades, embora as diferenças tenham sido pequenas em vários casos. O XGBoost permaneceu competitivo e foi superior em três localidades, especialmente nas séries com maior regularidade, enquanto a MLP venceu em duas. Na escala mensal, RNN, MLP e XGBoost alternaram a liderança entre as localidades. Essa diferença indica que não há um único modelo dominante para todas as regiões avaliadas. Portanto, a escolha do modelo deve considerar as características locais da série temporal, a escala temporal adotada e o comportamento observado no período de validação.

As métricas em W/m$^2$ facilitam a leitura física do erro. O RMSE foi consistentemente maior que o MAE, como esperado, pois penaliza mais fortemente erros pontuais de maior magnitude. O nRMSE permitiu comparar localidades com níveis médios de irradiância distintos: erros absolutos semelhantes podem representar impactos percentuais diferentes quando a média local de GHI é menor. Assim, localidades como GM Factory Zero e Ford Rouge Electric Vehicle Center apresentaram nRMSE diário elevado mesmo com R$^2$ ainda positivo, enquanto regiões como Lucid AMP 1 Casa Grande, Tesla Fremont Factory e Tesla Gigafactory Nevada apresentaram melhor combinação entre erro absoluto e poder explicativo.

Outro ponto importante é que o fluxo atual usa somente o histórico do próprio GHI. Dessa forma, diferenças não explicadas por defasagens e médias móveis, como variações rápidas de nebulosidade, umidade, aerossóis e outros fatores meteorológicos, aparecem como erro residual. Essa limitação é coerente com o escopo do trabalho e indica caminhos de continuidade para melhorar a capacidade explicativa dos modelos.

\vspace{12pt}

\noindent\textbf{Modelos avançados e comparação conjunta}

Os três modelos avançados foram executados para todas as dez localidades nas duas escalas temporais, totalizando 30 treinamentos diários e 30 treinamentos mensais concluídos com sucesso. A tabela a seguir resume as métricas médias por modelo avançado.

\vspace{6pt}

\begin{center}
\small
\begin{tabular}{llrrrr}
\toprule
\textbf{Freq.} & \textbf{Modelo} & \textbf{MAE} & \textbf{RMSE} & \textbf{R$^2$} & \textbf{nRMSE} \\
 & & \textbf{W/m$^2$} & \textbf{W/m$^2$} & & \textbf{\%} \\
\midrule
Diária & DilatedRNN & 37,70 & 49,71 & 0,656 & 26,86 \\
Diária & DeepNPTS\_aprox & 39,32 & 51,65 & 0,629 & 27,86 \\
Diária & DeepAR\_exp & 41,44 & 54,40 & 0,586 & 29,33 \\
Mensal & DeepNPTS\_aprox & 12,96 & 16,39 & 0,916 & 8,10 \\
Mensal & DilatedRNN & 15,25 & 18,73 & 0,878 & 9,29 \\
Mensal & DeepAR\_exp & 61,02 & 68,57 & 0,010 & 34,21 \\
\bottomrule
\end{tabular}
\end{center}

\vspace{6pt}

Na escala diária, a DilatedRNN apresentou RMSE médio de 49,71 W/m$^2$, valor praticamente equivalente ao RMSE médio dos melhores modelos de referência por localidade. Isso indica que a arquitetura dilatada conseguiu aproveitar a sequência curta de defasagens sem exigir variáveis meteorológicas adicionais. O DeepNPTS\_aprox também manteve desempenho competitivo, enquanto o DeepAR\_exp apresentou maior erro médio nesta configuração.

Na escala mensal, o melhor resultado médio entre os modelos avançados foi obtido pelo DeepNPTS\_aprox, com RMSE de 16,39 W/m$^2$, R$^2$ de 0,916 e nRMSE de 8,10\%. Esse resultado foi ligeiramente superior à média dos melhores modelos de referência mensais. A DilatedRNN também permaneceu competitiva, enquanto o DeepAR\_exp teve desempenho inadequado para a série mensal neste formato, sugerindo sensibilidade à quantidade reduzida de amostras mensais disponíveis.

\vspace{6pt}

\begin{center}
\small
\begin{tabular}{llr}
\toprule
\textbf{Freq.} & \textbf{Modelo ou grupo vencedor} & \textbf{Vitórias} \\
\midrule
Diária & Avançado - DilatedRNN & 5 \\
Diária & Referência - MLP & 2 \\
Diária & Referência - RNN & 2 \\
Diária & Referência - XGBoost & 1 \\
Mensal & Avançado - DeepNPTS\_aprox & 5 \\
Mensal & Avançado - DilatedRNN & 1 \\
Mensal & Referência - RNN & 3 \\
Mensal & Referência - XGBoost & 1 \\
\bottomrule
\end{tabular}
\end{center}

\vspace{6pt}

Considerando o maior R$^2$ em W/m$^2$ como critério de vitória, os modelos avançados venceram cinco das dez localidades na escala diária e seis das dez na escala mensal. Esse resultado mostra que as arquiteturas avançadas são relevantes dentro da comparação conjunta, mas também evidencia que os modelos de referência permanecem competitivos, especialmente por serem mais simples, mais fáceis de interpretar e menos dependentes de ajustes específicos.

\vspace{12pt}

\noindent\textbf{4. CONCLUSÕES}

Este trabalho apresentou um fluxo de aprendizado de máquina para previsão da GHI média diária do dia seguinte e da GHI média mensal do mês seguinte em dez localidades associadas a fábricas de veículos elétricos. A metodologia utilizou dados oficiais da NSRDB/NREL entre 2019 e 2024, com validação, agregação temporal, quantização, normalização, criação de variáveis temporais e avaliação cronológica.

Os resultados mostram que a previsão de um dia à frente é viável usando apenas histórico diário de GHI. A RNN apresentou melhor desempenho em cinco localidades, o XGBoost em três e a MLP em duas, enquanto a LSTM não foi vencedora nesta configuração experimental. Considerando os melhores modelos por localidade, obteve-se RMSE médio de 49,72 W/m$^2$ e R$^2$ médio de 0,66. No fluxo mensal, o RMSE médio dos melhores modelos caiu para 16,77 W/m$^2$ e o R$^2$ médio aumentou para 0,89, indicando que a escala mensal é mais previsível por suavizar variações diárias de curto prazo.

Na comparação conjunta, a DilatedRNN se destacou na previsão diária, vencendo cinco localidades, enquanto o DeepNPTS\_aprox foi especialmente competitivo na escala mensal, com o menor RMSE médio entre os modelos avaliados nesse fluxo. Por outro lado, o DeepAR\_exp não apresentou desempenho competitivo nesta configuração, principalmente no experimento mensal. Assim, os resultados indicam que a escolha do modelo deve considerar simultaneamente a escala temporal, a localidade e o equilíbrio entre desempenho, estabilidade e interpretabilidade.

Como trabalhos futuros, recomenda-se incluir variáveis meteorológicas adicionais, como temperatura, nebulosidade, umidade e velocidade do vento, além de avaliar horizontes de previsão maiores e modelos globais treinados com múltiplas localidades. A metodologia também pode ser reutilizada em estudos preliminares de dimensionamento e comparação de sistemas solares fotovoltaicos em ambientes industriais.

\vspace{12pt}

\noindent\textbf{\textit{Agradecimentos}}

Os autores agradecem à Universidade Estadual de Campinas (UNICAMP) e à Faculdade de Engenharia Elétrica e de Computação (FEEC) pelo apoio institucional ao desenvolvimento deste projeto.

\vspace{12pt}

\noindent\textbf{REFERÊNCIAS}

\vspace{6pt}

\noindent Alexandrov, A.; Benidis, K.; Bohlke-Schneider, M.; Flunkert, V.; Gasthaus, J.; Januschowski, T.; Maddix, D. C.; Rangapuram, S.; Salinas, D.; Schulz, J.; Stella, L.; Türkmen, A. C. e Wang, Y. (2019), ``GluonTS: Probabilistic time series models in Python'', \textit{arXiv preprint arXiv:1906.05264}.

\vspace{6pt}

\noindent Chang, S.; Zhang, Y.; Han, W.; Yu, M.; Guo, X.; Tan, W.; Cui, X.; Witbrock, M.; Hasegawa-Johnson, M. e Huang, T. S. (2017), ``Dilated recurrent neural networks'', \textit{Advances in Neural Information Processing Systems}.

\vspace{6pt}

\noindent Chen, T. e Guestrin, C. (2016), ``XGBoost: A scalable tree boosting system'', \textit{Proceedings of the 22nd ACM SIGKDD International Conference on Knowledge Discovery and Data Mining}, San Francisco, p. 785--794.

\vspace{6pt}

\noindent Géron, A. (2022), \textit{Hands-On Machine Learning with Scikit-Learn, Keras, and TensorFlow}, 3a ed., O'Reilly Media, Sebastopol.

\vspace{6pt}

\noindent Haykin, S. (2009), \textit{Neural Networks and Learning Machines}, 3a ed., Pearson, New Jersey.

\vspace{6pt}

\noindent Hyndman, R. J. e Athanasopoulos, G. (2021), \textit{Forecasting: Principles and Practice}, 3a ed., OTexts, Melbourne.

\vspace{6pt}

\noindent Pinho, J. T. e Galdino, M. A. (2014), \textit{Manual de Engenharia para Sistemas Fotovoltaicos}, CEPEL/CRESESB, Rio de Janeiro.

\vspace{6pt}

\noindent PVLIB Python Development Team (2026), \textit{pvlib python: a Python package for modeling solar energy systems}, documentação técnica.

\vspace{6pt}

\noindent Salinas, D.; Flunkert, V. e Gasthaus, J. (2017), ``DeepAR: Probabilistic forecasting with autoregressive recurrent networks'', \textit{arXiv preprint arXiv:1704.04110}.

\vspace{6pt}

\noindent Sengupta, M.; Xie, Y.; Lopez, A.; Habte, A.; Maclaurin, G. e Shelby, J. (2018), ``The National Solar Radiation Data Base (NSRDB)'', \textit{Renewable and Sustainable Energy Reviews}, vol. 89, p. 51--60.

\vspace{6pt}

\noindent Villalva, M. G. e Gazoli, J. R. (2012), \textit{Energia Solar Fotovoltaica: Conceitos e Aplicações}, 2a ed., Érica, São Paulo.

\vspace{6pt}

\noindent Wilcox, S. (2012), \textit{National Solar Radiation Database 1991--2010 Update: User's Manual}, National Renewable Energy Laboratory, Golden.

\end{document}
